%----------------------------------
\section{Data Set}
\label{sec:Data Set}
%----------------------------------

%%----------------------------------
\subsection{Description of the Data Sets}
\label{subsec:Description of the Data Sets}
%----------------------------------

Two public, account-free corpora of crowd-labelled user text underpin the analysis
(Table~\ref{tab:datasets}). Davidson et al.'s hate speech corpus~\cite{davidson2017automated}
holds 24,783 tweets, each judged by at least three CrowdFlower annotators and assigned one
of three classes: 5.77\% hate speech, 77.43\% offensive language, 16.80\% neither.
Wulczyn et al.'s Wikipedia Talk corpus~\cite{wulczyn2017ex} holds 159,686 discussion
comments carrying 1,598,289 toxicity judgements from 4,301 annotators, roughly ten per
comment; 9.62\% are toxic by majority vote. A companion file records each annotator's
demographics. Both download without an account from the URLs in Table~\ref{tab:datasets}.

\wordcount{90}

\begin{table}[h]
\centering
\caption{The two public corpora used in this case study. Class balance is computed from
the raw label files; the positive class is the item a moderator would be asked to act on.}
\label{tab:datasets}
\begin{tabular}{l>{\raggedright\arraybackslash}p{4.6cm}>{\raggedright\arraybackslash}p{4.6cm}}
\toprule
 & \textbf{Dataset A: Davidson et al.} & \textbf{Dataset B: Wulczyn et al.} \\
\midrule
Source & Twitter & English Wikipedia talk pages \\
URL & \url{https://github.com/t-davidson/hate-speech-and-offensive-language} &
      \url{https://figshare.com/articles/dataset/Wikipedia_Talk_Labels_Toxicity/4563973} \\
Instances & 24,783 tweets & 159,686 comments \\
File size & 2.4~MB (CSV) & 109~MB (two TSV files) \\
Attributes & \texttt{count}, \texttt{hate\_speech}, \texttt{offensive\_language},
             \texttt{neither}, \texttt{class}, \texttt{tweet} &
             \texttt{rev\_id}, \texttt{comment}, \texttt{year}, \texttt{logged\_in},
             \texttt{ns}, \texttt{sample}, \texttt{split} \\
Annotations & $\geq 3$ per tweet (median 3, max 9) & 1,598,289 judgements, 4,301
              annotators, $\approx$10 per comment \\
Label scheme & 3 classes: hate speech / offensive / neither & Binary toxicity per
               annotator, aggregated by majority \\
Positive class & Hate speech (5.77\%, $n=1{,}430$) & Toxic (9.62\%, $n=15{,}362$) \\
Mean length & 85.4 characters & 446.9 characters \\
Licence & Public research release (MIT) & CC0 1.0 Public Domain \\
\bottomrule
\end{tabular}
\end{table}

%%----------------------------------
\subsection{Rationale for Selection}
\label{subsec:Rationale for Selection}
%----------------------------------

Both corpora are exactly the modality the advertised role's classifiers consume: short,
user-generated, crowd-labelled text. Their label schemes deliberately disagree ---
Davidson separates hate speech from merely offensive language, while Wikipedia collapses
both into ``toxic'' --- which lets me test whether a triage model survives a change of
platform, the failure mode that matters in production. Limitations are real: both are
English-only, Davidson's lexicon-seeded sampling over-represents slurs, Wikipedia talk
pages are not video comments, neither includes audio or vision, and majority-vote labels
discard annotator disagreement.

\wordcount{83}
